% =====================================================================
% Fundamentos de Sistemas Electrónicos Analógicos — Resistores
% Compilar: pdflatex clase_2_resistores.tex  (2 pasadas)
% =====================================================================
\documentclass[aspectratio=169]{beamer}

\usepackage[utf8]{inputenc}
\usepackage[T1]{fontenc}
\usepackage{lmodern}
\usepackage[spanish,es-tabla]{babel}
\usepackage{amsmath,amssymb}
\usepackage{siunitx}
\usepackage{booktabs}
\usepackage{microtype}
\usepackage{circuitikz}

\definecolor{navy}{RGB}{23,54,93}
\definecolor{navysoft}{RGB}{72,101,140}
\definecolor{gold}{RGB}{179,142,62}
\definecolor{ink}{RGB}{44,51,58}
\definecolor{soft}{RGB}{244,241,234}

\usetheme{default}
\setbeamertemplate{navigation symbols}{}
\setbeamertemplate{footline}[frame number]
\setbeamercolor{structure}{fg=navy}
\setbeamercolor{frametitle}{fg=white,bg=navy}
\setbeamercolor{title}{fg=navy}
\setbeamercolor{subtitle}{fg=navysoft}
\setbeamercolor{itemize item}{fg=gold}
\setbeamercolor{itemize subitem}{fg=gold}
\setbeamercolor{block title}{fg=white,bg=navy}
\setbeamercolor{block body}{bg=soft}
\setbeamercolor{example text}{fg=navy}
\setbeamertemplate{itemize item}{\scriptsize\raise1.25pt\hbox{\donotcoloroutermaths$\bullet$}}
\setbeamertemplate{itemize subitem}{\scriptsize\raise1.25pt\hbox{\donotcoloroutermaths$\circ$}}
\setbeamerfont{frametitle}{size=\large,series=\bfseries}

\title[Resistores]{Resistores: fundamentos y aplicaciones}
\subtitle{Clase de teoría · Fundamentos de Sistemas Electrónicos Analógicos}
\institute{Licenciatura en Ingeniería Aeroespacial · UNAM}
\date{Semestre 2027-1}

\begin{document}

% ---------------------------------------------------------------------
\begin{frame}[plain]
  \titlepage
\end{frame}

% ---------------------------------------------------------------------
\begin{frame}{Objetivos}
  \begin{itemize}
    \item Definir un \textbf{resistor} y su papel en el circuito.
    \item Clasificar en \textbf{fijos, variables y no lineales}.
    \item Interpretar \textbf{tolerancia}, \textbf{series E} y \textbf{código de colores} (4/5/6 bandas).
    \item Modelar el \textbf{TCR} y la \textbf{potencia disipada}.
    \item Elegir el tipo constructivo y el valor para una aplicación.
  \end{itemize}
  \pause
  \begin{block}{Conexión con la semana}
    Todo lo de aquí se usa para \textbf{dimensionar} el divisor, el puente y el
    acondicionamiento de señales de las semanas 1–2.
  \end{block}
\end{frame}

% ---------------------------------------------------------------------
\begin{frame}{1 · ¿Qué es un resistor?}
  \begin{columns}[T]
    \begin{column}{0.45\textwidth}
      \begin{center}
        \resizebox{0.9\linewidth}{!}{%
        \begin{circuitikz}[american voltages, scale=1.0]
          \draw (0,2.6) -- (2.6,2.6) to[R, l={$R$}] (5.2,2.6) -- (6.6,2.6);
          \node[left] at (0,2.6) {$+$};
        \end{circuitikz}}
      \end{center}
      \begin{itemize}
        \item Un resistor \textbf{se opone al paso de la corriente}.
        \item Transforma energía eléctrica en \textbf{calor}.
      \end{itemize}
    \end{column}
    \begin{column}{0.5\textwidth}
      \[
        V = I\cdot R
        \qquad
        I = \frac{V}{R}
      \]
      \textbf{Unidad:} ohm ($\Omega$). \textbf{Ley de Ohm:} la corriente es
      proporcional a la tensión e inversa a la resistencia.
      \pause
      \begin{block}{Punto de trabajo}
        Todo componente tiene un punto de reposo; el resistor lo fija/limita.
        Sin resistor, corrientes no controladas y sobrecalentamiento.
      \end{block}
    \end{column}
  \end{columns}
\end{frame}

% ---------------------------------------------------------------------
\begin{frame}{2 · Clasificación de los resistores}
  \begin{table}
    \small
    \begin{tabular}{p{0.26\textwidth}p{0.34\textwidth}p{0.32\textwidth}}
      \toprule
      Tipo & Característica & Ejemplo \\
      \midrule
      \textbf{Fijos} & $R$ constante (solo cambia con T/tolerancia) & divisor, limitación de LED \\
      \textbf{Variables} & $R$ ajustable con un cursor (potenciómetro) & control de volumen, divisor ajustable \\
      \textbf{No lineales} & $R$ cambia con una magnitud física & termistor, LDR, varistor \\
      \bottomrule
    \end{tabular}
  \end{table}
  \pause
  \begin{itemize}
    \item Un resistor \textbf{ideal} sería constante; el \textbf{real} depende de
      tolerancia, temperatura y (si es no lineal) de la magnitud a sensar.
  \end{itemize}
\end{frame}

% ---------------------------------------------------------------------
\begin{frame}{3 · Fijos: valor nominal y tolerancia}
  El valor real puede desviarse del nominal dentro de la \textbf{tolerancia}:
  \[
    R_{min} = R\left(1-\frac{t}{100}\right)
    \qquad
    R_{max} = R\left(1+\frac{t}{100}\right)
  \]
  \textbf{Ejemplo:} $1\,\text{k}\Omega \pm 5\text{\%}$ $\Rightarrow$ \SI{950}{\ohm} a \SI{1050}{\ohm}.
  \pause
  \begin{itemize}
    \item Tolerancia típica: $\pm5\text{\%}$ (E24), $\pm1\text{\%}$ (E96), $\pm0.1\text{\%}$ (precisión).
    \item En un divisor, importa la \emph{razón} (ver S2.2): los cambios comunes se cancelan.
  \end{itemize}
\end{frame}

% ---------------------------------------------------------------------
\begin{frame}{4 · Fijos: series normalizadas (IEC 60063)}
  No existe "cualquier valor": se usa el de la \textbf{serie} más cercano al cálculo.
  \begin{itemize}
    \item \textbf{E12} ($\pm10\text{\%}$): \texttt{1.0 1.2 1.5 1.8 2.2 2.7 3.3 3.9 4.7 5.6 6.8 8.2}
    \item \textbf{E24} ($\pm5\text{\%}$): \texttt{1.0 1.1 1.2 1.3 1.5 1.6 1.8 2.0 2.2 2.4 2.7 3.0 3.3 3.6 3.9 4.3 4.7 5.1 5.6 6.2 6.8 7.5 8.2 9.1}
    \item \textbf{E96} ($\pm1\text{\%}$): 3 cifras (más fina).
  \end{itemize}
  \pause
  \textbf{Ejemplo:} cálculo da $2.37\,\text{k}\Omega$ $\Rightarrow$ E24: \texttt{2.4 k} (error 1.3\text{\%}); E96: \texttt{2.37 k} (error 0\text{\%}).
\end{frame}

% ---------------------------------------------------------------------
\begin{frame}{5 · Fijos: código de colores (IEC 60062)}
  Se leen de izquierda a derecha: \textbf{4} (dígito, dígito, mult., tol.), \textbf{5}
  (3 dígitos), \textbf{6} (agrega una banda de \textbf{TCR}).
  \begin{center}
    \resizebox{0.92\linewidth}{!}{%
    \begin{tabular}{llll}
      \toprule
      Color & Dígito & Multiplicador & Tolerancia \\
      \midrule
      Negro & 0 & $\times1$ & — \\
      Marrón & 1 & $\times10$ & $\pm1\text{\%}$ \\
      Rojo & 2 & $\times100$ & $\pm2\text{\%}$ \\
      Naranja & 3 & $\times1$k & — \\
      Amarillo & 4 & $\times10$k & — \\
      Verde & 5 & $\times100$k & $\pm0.5\text{\%}$ \\
      Azul & 6 & $\times1$M & $\pm0.25\text{\%}$ \\
      Violeta & 7 & — & $\pm0.1\text{\%}$ \\
      Gris & 8 & — & $\pm0.05\text{\%}$ \\
      Blanco & 9 & — & — \\
      Dorado & — & $\times0.1$ & $\pm5\text{\%}$ \\
      Plata & — & $\times0.01$ & $\pm10\text{\%}$ \\
      \bottomrule
    \end{tabular}}
  \end{center}
  \textbf{Ejemplo:} \texttt{Marrón–Negro–Rojo–Dorado} $\Rightarrow$ $10\times100 = \SI{1}{\kilo\ohm}\ \pm5\text{\%}$.
\end{frame}

% ---------------------------------------------------------------------
\begin{frame}{6 · Fijos: coeficiente de temperatura (TCR)}
  El valor cambia con la temperatura (\textbf{TCR}, ppm/°C):
  \[
    \Delta R = R_{nominal}\cdot TCR\cdot \Delta T
    \qquad
    error\text{\%} = \frac{\Delta R}{R}
  \]
  \textbf{Ejemplo:} $10\,\text{k}\Omega$, $100\,\text{ppm/°C}$, $\Delta T = 50\,^\circ\text{C}$
  $\Rightarrow \Delta R = \SI{50}{\ohm}$ (0.5\text{\%}).
  \pause
  \begin{itemize}
    \item \textbf{Carbón}: $\sim100$–$200$ ppm/°C. \textbf{Película metálica}: $\le 50$ ppm/°C.
    \item En un divisor de precisión (referencia de ADC) se exige TCR \textbf{bajo}.
  \end{itemize}
\end{frame}

% ---------------------------------------------------------------------
\begin{frame}{7 · Fijos: potencia nominal y disipación}
  \[
    P = V\cdot I = I^2\,R = \frac{V^2}{R}
  \]
  El resistor debe soportar la potencia disipada; \textbf{en el vacío} (sin convección)
  las normas aplican \textbf{derating} ($\approx 50\text{\%}$):
  \[
    P_{derated} = 0.5\cdot P_{nominal}
  \]
  \textbf{Ejemplo:} $1\,\text{k}\Omega$ de $1/4\,\text{W}$ a \SI{12}{\volt}
  $\Rightarrow P = \SI{144}{\milli\watt} > \SI{125}{\milli\watt}$ (límite) $\Rightarrow$ usar $1/2\,\text{W}$.
\end{frame}

% ---------------------------------------------------------------------
\begin{frame}{8 · Fijos: tipos constructivos}
  \begin{table}
    \small
    \begin{tabular}{p{0.24\textwidth}p{0.30\textwidth}p{0.30\textwidth}}
      \toprule
      Tipo & Precisión/TCR & Ruido y uso \\
      \midrule
      Composición de carbono & $\pm10\text{\%}$, TCR alto & Alto ruido (evitar en aeroespacial) \\
      Película de carbono & $\pm5\text{\%}$, TCR medio & Uso general \\
      Película metálica & $\pm0.1$–$1\text{\%}$, TCR bajo & \textbf{Señales/sensores} \\
      Bobinado & Buena precisión, alta potencia & Inductivo (no para RF) \\
      SMD (montaje superficial) & Excelente & \textbf{Vibración/avionica} \\
      \bottomrule
    \end{tabular}
  \end{table}
  \pause
  \begin{itemize}
    \item \textbf{En vuelo} se prefiere \textbf{SMD} y \textbf{película metálica}:
      baja masa (menos fatiga en vibración) y bajo ruido.
  \end{itemize}
\end{frame}

% ---------------------------------------------------------------------
\begin{frame}{9 · Variables: potenciómetro y reóstato}
  \begin{columns}[T]
    \begin{column}{0.46\textwidth}
      \resizebox{\linewidth}{!}{%
      \begin{circuitikz}[american voltages, scale=0.95]
        \draw (0,0) to[V, l={$V_{in}$}] (0,2.6);
        \draw (0,2.6) -- (1.4,2.6) to[R, l={$R_{top}$}] (4.4,2.6);
        \draw (4.4,2.6) -- (5.4,2.6) to[R, l={$R_{bot}$}] (5.4,0);
        \draw (5.4,0) -- (0,0);
        \node[circle,fill=gold,inner sep=0.9pt] at (4.4,2.6) {};
        \node[above,text=gold] at (4.4,2.9) {$V_{out}$};
      \end{circuitikz}}
    \end{column}
    \begin{column}{0.5\textwidth}
      \begin{itemize}
        \item \textbf{Potenciómetro}: divisor con cursor; $V_{out} = V_{in}\,[R_{bot}/(R_{top}+R_{bot})]$.
        \item \textbf{Reóstato}: solo dos terminales (resistencia variable en serie).
        \item \textbf{Trimpot}: ajuste de precisión con destornillador.
        \item \textbf{Lineal vs logarítmico}: la curva puede ser lineal o de audio.
      \end{itemize}
    \end{column}
  \end{columns}
\end{frame}

% ---------------------------------------------------------------------
\begin{frame}{10 · No lineales: sensores y protección}
  \begin{table}
    \small
    \begin{tabular}{p{0.22\textwidth}p{0.40\textwidth}p{0.30\textwidth}}
      \toprule
      Tipo & $R$ varía con & Uso \\
      \midrule
      \textbf{NTC} & temperatura (baja al subir $T$) & sensor/sensor de temperatura \\
      \textbf{PTC} & temperatura (sube al subir $T$) & protección por sobrecorriente \\
      \textbf{RTD (Pt100)} & temperatura (lineal, precisa) & medición de precisión \\
      \textbf{Varistor} & tensión (cae ante sobretensión) & \textbf{protección} contra transitorios \\
      \textbf{LDR} & luz (baja al iluminar) & detección de luz \\
      \bottomrule
    \end{tabular}
  \end{table}
  \pause
  \begin{itemize}
    \item En aeroespacial, \textbf{RTD} (Pt100) para temperatura precisa; \textbf{varistor}
      para proteger el bus de transitorios; un \textbf{termistor NTC} para sensor barato.
  \end{itemize}
\end{frame}

% ---------------------------------------------------------------------
\begin{frame}{Aplicación aeroespacial}
  \begin{itemize}
    \item \textbf{Precisión:} divisores de referencia requieren \textbf{película metálica}
      (tol. $\pm0.1\text{\%}$, TCR $\le 25$ ppm/°C).
    \item \textbf{Robustez:} en el \textbf{vacío} no hay convección $\Rightarrow$ \textbf{derating}
      y acoplamiento térmico al chasis; el \textbf{SMD} resiste vibración.
    \item \textbf{Interfaz:} el acondicionamiento usa divisores, puentes de Wheatstone
      (con RTD/galga) y \textbf{equivalente de Thévenin} para las etapas de amplificación y ADC.
    \item \textbf{Selección:} se elige el resistor con \emph{valor} + \emph{tolerancia} +
      \emph{TCR} + \emph{potencia con derating}, y se documenta en la \emph{ficha} de la interfaz.
  \end{itemize}
\end{frame}

% ---------------------------------------------------------------------
\begin{frame}{Resumen}
  \begin{itemize}
    \item El resistor \textbf{opone} paso a corriente y \textbf{disipa} calor ($V=IR$, $P=V^2/R$).
    \item Se clasifica en \textbf{fijo, variable y no lineal}.
    \item El fijo se elige por \textbf{valor} (serie E), \textbf{tolerancia}, \textbf{TCR} y \textbf{potencia (derating)}.
    \item \textbf{Película metálica + SMD} para aeroespacial; \textbf{RTD} para temperatura precisa.
    \item El divisor (con carga) y su equivalente conectan con el acondicionamiento de señales.
  \end{itemize}
  \pause
  \begin{block}{Frase clave}
    Un resistor "que funciona" en el cálculo no basta: hay que elegirlo con
    \textbf{tolerancia}, \textbf{deriva térmica} y \textbf{potencia} para la misión.
  \end{block}
\end{frame}

% ---------------------------------------------------------------------
\begin{frame}{Preguntas de verificación}
  \begin{enumerate}
    \item ¿Qué describe la tolerancia y en qué unidades se da el TCR?
      \uncover<2->{\emph{(El desvío del valor (\text{\%}); TCR en ppm/°C.)}}
    \item Lee \texttt{Amarillo–Violeta–Naranja–Plata}: ¿valor y tolerancia?
      \uncover<3->{\emph{($47\times1\text{k} = \SI{47}{\kilo\ohm}$, $\pm10\text{\%}$.)}}
    \item ¿Cuál usarías para un divisor de precisión: carbón o película metálica? ¿Por qué?
      \uncover<4->{\emph{(Película metálica: menor TCR y tolerancia.)}}
    \item ¿Qué tipo de resistor cambia su valor con la luz? ¿Y con la tensión?
      \uncover<5->{\emph{(LDR; varistor.)}}
    \item ¿Por qué en el vacío se debe derating de potencia?
      \uncover<6->{\emph{(No hay convección; el sobrecalentamiento se agrava.)}}
  \end{enumerate}
\end{frame}

\end{document}
